\documentclass[11pt]{article}
\usepackage[margin=25mm]{geometry}
\usepackage{fontspec}
\setmainfont{Latin Modern Roman}
\usepackage{amsmath,amssymb,mathtools}
\usepackage{unicode-math}
\setmathfont{Latin Modern Math}
\usepackage{xurl,hyperref}
\hypersetup{colorlinks=true,urlcolor=blue}
\setcounter{secnumdepth}{0}
\setlength{\parindent}{0pt}
\setlength{\parskip}{5pt}
\setlength{\emergencystretch}{3em}
\newcommand{\sref}[2]{\hyperlink{source:#1:#2}{[S#2]}}
\newcommand{\cataloguescope}[1]{\paragraph{Recorded target.}#1}
\begin{document}
\section{The Hodge Conjecture}\label{statement}

\subsection{Definitions and mathematical statement}

\paragraph{Shared notation.}$\mathbb N=\{1,2,\ldots\}$, and $\mathbb Z,\mathbb Q,\mathbb R,\mathbb C$ denote the integers, rationals, reals and complex numbers. Any other conventions are specified locally.


Let \(X\) be a smooth projective variety over \(\mathbb C\).
Its singular cohomology has Hodge decomposition
\[
 H^k(X,\mathbb C)=\bigoplus_{a+b=k}H^{a,b}(X).
\]
A rational Hodge class of codimension \(p\) is a class
\(\alpha\in H^{2p}(X,\mathbb Q)\) whose image in complex cohomology
lies in \(H^{p,p}(X)\).  For a closed irreducible algebraic subvariety
\(Z\subseteq X\) of complex codimension \(p\), write
\([Z]\in H^{2p}(X,\mathbb Q)\) for its cycle class, Poincar\'e dual to
its fundamental homology class.  A rational algebraic cycle is a finite
formal sum of such subvarieties with rational coefficients.
\paragraph{Statement recorded in the supplied source.}
For every smooth complex projective variety \(X\), every integer
\(0\le p\le\dim_{\mathbb C}X\), and every rational Hodge class
\(\alpha\in H^{2p}(X,\mathbb Q)\), do there exist algebraic subvarieties
\(Z_1,\ldots,Z_r\subseteq X\) of codimension \(p\) and numbers
\(q_1,\ldots,q_r\in\mathbb Q\) such that
\[
 \alpha=\sum_{i=1}^{r}q_i[Z_i]?
\] \sref{statement}{2}
\subsection{Short English statement}
Is every rational cohomology class of Hodge type (p,p) on a smooth complex projective variety represented by a rational algebraic cycle?
\subsection{Sources}
\begin{description}
\item[\hypertarget{source:statement:1}{S1}] ULAM.AI and the UnsolvedMath Contributors, \emph{UnsolvedMath: A Curated Collection of Open Mathematics Problems}, catalogue entry. CC BY 4.0. \url{https://huggingface.co/datasets/ulamai/UnsolvedMath}.
\item[\hypertarget{source:statement:2}{S2}] Claire Voisin, \emph{A counterexample to the Hodge conjecture for K\"ahler varieties}, arXiv:math/0112247 (2001), Section 1, page 1, statement of the projective Hodge conjecture. \url{https://arxiv.org/abs/math/0112247}.
\end{description}
\label{end:statement}
\end{document}
